% !PRINTED 693
\weberpage{693}
%<<CONT
und $\varrho$ könnte nicht durch $r\varepsilon$ teilbar sein. Ebenso schließt man,
daß jeder der Primfaktoren $\pi_1, \pi_2, \ldots$, aber keiner mehr als einmal
 in $r\varepsilon$ aufgeht; also ist, von Einheitsfaktoren abgesehen,

\weq{15}{r\varepsilon = \varkappa.}

Also genügt jede der Funktionen $\varepsilon_1, \varepsilon_2, \ldots, \varepsilon_n$ der Bedingung a).
Umgekehrt ist auch jede dieser Bedingung genügende Funktion $w$
in der Form enthalten:

\weq{16}{w = x_1\varepsilon_1 + x_2\varepsilon_2 + \cdots + x_n\varepsilon_n,}

worin die $x_1, x_2, \ldots, x_n$ ganze rationale Funktionen von $z$ sind.

Denn es ist $rw$ eine durch $\varkappa$ teilbare ganze Funktion, und
folglich nach (9), (16) und § 184, 7.:
\[
S(rw\omega_i) = rx_i
\]
eine durch $r$ teilbare ganze Funktion von $z$, folglich

\weq{17}{x_i = S(w\omega_i).}

Um also die Differentialquotienten erster Gattung zu erhalten,
hat man unter den Funktionen (16) die auszusuchen, die
der Bedingung b) genügen.

Zu dem Zweck nehmen wir für (2) die in § 193 betrachtete
Normalbasis
\[
\lambda_1, \lambda_2, \ldots, \lambda_n
\]
und bezeichnen die dazu komplementäre Basis mit
\[
\mu_1, \mu_2, \ldots, \mu_n,
\]
die durch

\weq{18}{S(\lambda_r\mu_s) = (r, s)}

definiert ist.

Setzen wir dann

\weq{19}{w = y_1\mu_1 + y_2\mu_2 + \cdots + y_n\mu_n,}

so ist, wenn $r_i$ der Exponent von $\lambda_i$ ist, nach (17)
\[
\frac{y_i}{z^{r_i-1}} = S\!\left(wz\frac{\lambda_i}{z^{r_i}}\right)\!.
\]

Nach der Definition von $r_i$ ist $\lambda_i:z^{r_i}$ für $z=\infty$ endlich und
$wz$ verschwindet wegen b). Folglich muß $y_i:z^{r_i-1}$ für $z=\infty$
verschwinden, d.\,h. $y_i$ kann höchstens vom Grade $(r_i-2)$ sein
und enthält daher höchstens $r_i-1$ konstante Koeffizienten.

Da $\mu_1=1$, $r_1=0$ ist und der Grad von $y_1$ nicht negativ
% [sic] printed $\mu_1=1$, but the fact used is $\lambda_1=1$ ($r_1=0$ is the exponent of $\lambda_1$, and $y_1=S(w\lambda_1)=S(w)$ gives $S(w)=0$); with $\lambda_1=1$, (18) gives $S(\mu_1)=1\neq S(1)=n$, so $\mu_1\neq1$
sein kann, so muß $y_1$ identisch Null sein und es folgt
\[
S(w) = 0.
\]
